\clearpage
\chapter*{Section Annexes}
\addcontentsline{toc}{chapter}{Annexes}

\section*{Annexe 1 — Vue ``Risk Intrinsink (RI)'' du dashboard Power BI IPANEMA}
\addcontentsline{toc}{section}{Annexe 1 — Vue ``Risk Intrinsink (RI)'' du dashboard Power BI IPANEMA}

Cette vue restitue la répartition des Tierces Parties selon leur niveau de Risque
Intrinsèque (LOW, MEDIUM, HIGH), calculé conformément à la méthodologie décrite en
2.2.3. Un graphique en anneau présente la répartition globale des niveaux de risque,
complété par un histogramme croisant ces niveaux avec le type de Tierce Partie
(sous-traitant clinique ou non clinique, prestataire de service, CDMO, distributeur
commercial et/ou pharmaceutique). Un tableau détaillé liste, pour chaque partenaire, son
type de tiers, son activité, son type de sous-traitance et son niveau de RI, permettant une
lecture individuelle en complément de la vue synthétique.

\begin{figure}[H]
    \centering
    \roundedfullwidthimage{cover/figures/annex1.png}
    \label{fig:annexe1_risk_intrinsink}
\end{figure}

\section*{Annexe 2 — Vue ``Risk Compliance \& Performance (RCP) — Tous''\\du dashboard Power BI IPANEMA}
\addcontentsline{toc}{section}{Annexe 2 — Vue ``Risk Compliance \& Performance (RCP) — Tous'' du dashboard Power BI IPANEMA}

Cette vue présente la répartition de l'ensemble des Tierces Parties selon leur niveau
de Risque Conformité et Performance, sans distinction de statut CDMO. Comme pour la vue
RI, un graphique en anneau et un histogramme par type de tiers offrent une vision
synthétique, tandis qu'un tableau restitue le détail par partenaire, incluant la note RCP
obtenue et le niveau de risque correspondant.

\begin{figure}[H]
    \centering
    \roundedfullwidthimage{cover/figures/annex2.png}
    \label{fig:annexe2_rcp_tous}
\end{figure}

\section*{Annexe 3 — Vue ``Risk Compliance \& Performance (RCP) — CDMO''\\du dashboard Power BI IPANEMA}
\addcontentsline{toc}{section}{Annexe 3 — Vue ``Risk Compliance \& Performance (RCP) — CDMO'' du dashboard Power BI IPANEMA}

Cette vue isole les Tierces Parties de statut CDMO afin de restituer leur niveau de
RCP selon une logique de calcul spécifique (cf. 2.5.5, mesure RCP\_CDMO\_Level). Le
tableau associé présente, pour chaque partenaire, la moyenne des notes RCP obtenues, le
premier niveau de risque enregistré ainsi que la moyenne toutes campagnes confondues,
permettant de suivre l'évolution du niveau de risque d'une Tierce Partie CDMO dans le
temps.

\begin{figure}[H]
    \centering
    \roundedfullwidthimage{cover/figures/annex3.png}
    \label{fig:annexe3_rcp_cdmo}
\end{figure}

\section*{Annexe 4 — Vue ``Risk Global (RG) — Tous'' du dashboard Power BI IPANEMA}
\addcontentsline{toc}{section}{Annexe 4 — Vue ``Risk Global (RG) — Tous'' du dashboard Power BI IPANEMA}

Cette vue restitue la répartition de l'ensemble des Tierces Parties selon leur niveau
de Risque Global, résultant de la combinaison matricielle du RI et du RCP décrite en 2.2.3.
L'histogramme distingue les quatre niveaux possibles — LOW, MEDIUM, HIGH et VERY
HIGH — croisés avec le type de tiers, offrant une vision d'ensemble du niveau de risque final
porté par chaque catégorie de partenaires.

\begin{figure}[H]
    \centering
    \roundedfullwidthimage{cover/figures/annex4.png}
    \label{fig:annexe4_risk_global_tous}
\end{figure}

\section*{Annexe 5 — Vue ``TP Qualification Status'' du dashboard Power BI IPANEMA}
\addcontentsline{toc}{section}{Annexe 5 — Vue ``TP Qualification Status'' du dashboard Power BI IPANEMA}

Cette vue restitue le statut de qualification de chaque Tierce Partie tel qu'enregistré
dans TrackWise, indépendamment de la campagne d'évaluation en cours. Le tableau
détaille, pour chaque partenaire, son statut actuel (qualifiée ou en attente de qualification), le
commentaire d'audit associé ainsi que son niveau de Risque Intrinsèque. Cette vue a fait
l'objet d'une correction spécifique portant sur la sélection du dernier enregistrement en date
par Tierce Partie, afin de ne restituer que le statut réellement à jour (cf. 2.4.2).

\begin{figure}[H]
    \centering
    \roundedfullwidthimage{cover/figures/annex5.png}
    \label{fig:annexe5_tp_qualification_status}
\end{figure}

\section*{Annexe 6 — Vue ``Evaluated Projects List'' du dashboard Power BI IPANEMA}
\addcontentsline{toc}{section}{Annexe 6 — Vue ``Evaluated Projects List'' du dashboard Power BI IPANEMA}

Cette vue liste l'ensemble des projets évalués, rattachés à leur Tierce Partie
respective. Le tableau présente, pour chaque projet, sa référence commerciale, son code de
spécification, sa stratégie de développement, son étape d'avancement, ainsi que les
commentaires renseignés par l'utilisateur ayant créé l'évaluation.

\begin{figure}[H]
    \centering
    \roundedfullwidthimage{cover/figures/annex6.png}
    \label{fig:annexe6_evaluated_projects_list}
\end{figure}

\section*{Annexe 7 — Vue ``Quality Agreements'' du dashboard Power BI IPANEMA}
\addcontentsline{toc}{section}{Annexe 7 — Vue ``Quality Agreements'' du dashboard Power BI IPANEMA}

Cette vue recense, pour chaque Tierce Partie, l'existence et les caractéristiques de
son accord qualité enregistré dans TrackWise. Le tableau croise cette information avec le
niveau de Risque Intrinsèque du partenaire, son statut de qualification, le statut de son
produit ainsi que la date de mise à jour de l'accord, offrant une vision consolidée des
engagements qualité en vigueur.

\begin{figure}[H]
    \centering
    \roundedfullwidthimage{cover/figures/annex7.png}
    \label{fig:annexe7_quality_agreements}
\end{figure}

\section*{Annexe 8 — Vue ``Satisfaction'' du dashboard Power BI IPANEMA}
\addcontentsline{toc}{section}{Annexe 8 — Vue ``Satisfaction'' du dashboard Power BI IPANEMA}

Cette vue compare la satisfaction obtenue par chaque Tierce Partie sur les six
critères d'évaluation RCP Métier (compétences techniques, innovation, qualité des livrables,
réactivité et respect des délais, relationnel, traçabilité). Les histogrammes appliquent une
mise en forme conditionnelle par couleur, accompagnée d'un avertissement rappelant
explicitement le sens de lecture de l'échelle — une précaution destinée à éviter toute
confusion entre un niveau de satisfaction élevé et un niveau d'insatisfaction élevé.

\begin{figure}[H]
    \centering
    \roundedfullwidthimage{cover/figures/annex8.png}
    \label{fig:annexe8_satisfaction}
\end{figure}
    
